\section{An elementary support bound}\label{ms:support}

We now show that the contact inequality confines an extremal law to a
short interval. Because the predecessor term vanishes, it suffices to bound
the Gaussian part of the contact inequality at one radius and to show that
the resulting upper bound decreases along each exterior ray.

Throughout this section $n\ge42$, $(P,H)$ is an
extremal pair, and $C_n\ge\CE$. We also assume
$C_{n-1}\le C_n$, so the positive predecessor loss is zero.
Set $v=n-1$, $s=\sqrt{n/v}$ and $F_v(x)=\Phi(x/\sqrt v)$. For every $t$ in
the support of $P$, the contact equality of Lemma~\ref{ms:contact} gives
\begin{equation}\label{ms:04:contact}
 0\le H_n(H,t)+C_n\{-|t|^3+\tfrac32\rho t^2+3\beta t+\rho e_n\},
\end{equation}
where
\[
 e_n=n\bigl(\sqrt{n/(n-1)}-1\bigr)-\tfrac12<\frac{0.384}{n},
\]
and the exact Gaussian expression is
\[
 \frac{H_n(H,t)}{n^{3/2}}
 =F_v(H-t)-F_{v+1}(H)+tF_{v+1}'(H)
       -\frac{t^2-1}{2}F_{v+1}''(H).
\]
Here we used $R_{n-1}(P)\le C_{n-1}\le C_n$: the term involving the
predecessor constant is then nonpositive and has been dropped.
For the numerical bound on $e_n$, rationalization gives
\[
 ne_n=\frac{s^2(2s+1)}{2(s+1)^2}.
\]
This expression increases for $s\ge1$: its logarithmic derivative is
$2/s+2/(2s+1)-2/(s+1)>0$. At $s=81/80$ its value is
$793881/2073680<48/125$. Thus $e_n<0.384/n\le8/875$.

\subsection{Fourth-order expansion with uniform remainder}
Using $\partial_vF_v=F_v''/2$, expand $F_v(H-t)$ spatially through order
four, $F_{v+1}$ through second order in variance, and its first and second
spatial derivatives through first order in variance. The order-zero, linear
and quadratic spatial terms cancel. The remaining cubic and quartic
coefficients are $(-t^3+3t)/6$ and $(t^4-6t^2+3)/24$. Thus, for $t=\pm u$,
\begin{equation}\label{ms:04:expansion}
 \frac{H_n(H,\pm u)}{s^3}
  =\pm A(1-z_v^2)\phi(z_v)+B\phi'''(z_v)+R,\qquad z_v=H/\sqrt v,
\end{equation}
where
\[
 A=\frac{u^3-3u}{6},\qquad B=\frac{u^4-6u^2+3}{24\sqrt v}.
\]
For $u\ge5/2$, both are positive. The four Taylor remainders, after division
by $s^3$, have bounds
\[
 \frac{u^5\|\phi^{(4)}\|_\infty}{120v},\quad
 \frac{\|\phi^{(5)}\|_\infty}{48v^{3/2}},\quad
 \frac{u\|\phi^{(4)}\|_\infty}{8v},\quad
 \frac{(u^2-1)\|\phi^{(5)}\|_\infty}{16v^{3/2}}.
\]
Gaussian Fourier inversion gives
$\|\phi^{(4)}\|_\infty\le3\phi(0)<6/5$ and
$\|\phi^{(5)}\|_\infty\le8/\pi<51/20$. Therefore
\begin{equation}\label{ms:04:error}
 |R|\le E_v(u):=\frac{u^5/100+3u/20}{v}
       +\frac{51(3u^2-2)}{960v^{3/2}}.
\end{equation}
Every remainder bound is uniform in the raw threshold $H$.
Here we used $\pi>160/51$.

\begin{lemma}[Two all-threshold envelopes]\label{ms:04:envelopes}
For $u\ge5/2$,
\[
 H_n(H,\pm u)\le s^3\{\phi(0)(A+2B)+E_v(u)\}.
\]
If also $B\le A/4$, then
\[
 H_n(H,\pm u)\le
 s^3\left\{\phi(0)\left(A+\frac{9B^2}{4A}\right)+E_v(u)\right\}.
\]
Both assertions hold for every real $H$ and both signs separately.
\end{lemma}
\begin{proof}
The first follows from
$\|(1-z^2)\phi(z)\|_\infty=\phi(0)$ and
$\|\phi'''\|_\infty<2\phi(0)$.
For the latter bound, on $|z|\le\sqrt3$ the polynomial
$|3z-z^3|$ is at most $2$, with equality only at $|z|=1$ where the
Gaussian factor is smaller than one. Outside that interval,
$(|z|^3-3|z|)e^{-z^2/2}\le |z|^3e^{-z^2/2}
\le3\sqrt3e^{-3/2}<2$. The last inequality follows, for example,
from $e^{3/2}>1+3/2+9/8+27/48>4$ and $3\sqrt3<8$.
For the positive sign and $|z_v|\le1$, the retained terms in
\eqref{ms:04:expansion} are at most
$\phi(0)(A-Az_v^2+3B|z_v|)$; completing the square proves the sharper bound.
For $|z_v|>1$ the first term is nonpositive and the second is at most
$2B\phi(0)\le A\phi(0)/2$.
For the negative sign, its first term is at most
$2e^{-3/2}A\phi(0)<A\phi(0)/2$, and the second is at most
$2B\phi(0)\le A\phi(0)/2$. Add the uniform remainder.
\end{proof}

\subsection{Three support bounds}
The first radius uses only the global moment cap. Later contractions use
sharper moment bounds on probability panels. We state them in terms of the moment bounds they require.

Use
\[
 S=81/80,\quad  c_0=819/2000<\CE,\quad
 H_4=u^4-6u^2+3,
\]
so $s\le S$, $\sqrt v>6$ and $e_n<8/875$. Define
\begin{align*}
 E(u)&=\frac{u^5}{4100}+\frac{3u}{820}
                   +\frac{17(3u^2-2)}{78720},\\
 Q_{r,b}(u)&=-u^3+\tfrac32ru^2+3bu+\tfrac8{875}r,\\
 P_{r,b}(u)&=\tfrac25S^3(A+H_4/72)+S^3E(u)+c_0Q_{r,b}(u),\\
 P^\sharp_{r,b}(u)&=\tfrac25S^3
       \left(A^2+\frac{H_4^2}{256\cdot41}\right)
       +S^3E(u)A+c_0Q_{r,b}(u)A.
\end{align*}
If $\rho\le r$ and $|\beta|\le b$, the contact coefficient is at
most $Q_{r,b}$. Check $Q_{r,b}<0$ before replacing its actual multiplier
$C_n$ by the lower constant $c_0$. Lemma~\ref{ms:04:envelopes} then makes
$P<0$ sufficient at a boundary; $P^\sharp<0$ is sufficient when $B\le A/4$.

To extend a boundary to a whole ray, put
\[
 L_{r,b}(u)=\tfrac25(u^2/2+1/2+8/875)+c_0Q'_{r,b}(u).
\]
To bound the derivative of $H_n$, put
$z=H/\sqrt n$ and $y=(H-t)/\sqrt n$. Exact differentiation gives
\[
 \partial_tH_n=n\phi(z)+Ht\phi(z)-ns\phi(sy).
\]
Taylor's theorem and $\|\phi''\|_\infty=\phi(0)$ imply
$|n\phi(y)-n\phi(z)-Ht\phi(z)|\le\phi(0)t^2/2$.
Also $\partial_s[s\phi(sy)]=-\phi''(sy)$, whence
$|s\phi(sy)-\phi(y)|\le\phi(0)(s-1)$. Since
$n(s-1)=1/2+e_n$, this proves
\[
 |\partial_tH_n|\le\phi(0)(t^2/2+1/2+e_n)
                    <(2/5)(t^2/2+1/2+e_n).
\]
If $Q,Q',Q'',L,L',L''$ are negative at $U$, their required signs persist
for $u\ge U$, since $Q'''=-6$ and $L''=2/5-6c_0<0$.
The contact upper bound for each sign then decreases strictly along the
whole ray.

\begin{proposition}[Support from moment bounds]\label{ms:04:three-bounds}
Under the contact and zero-loss hypotheses of this section, each of the
following implies $\supp P\subset(-R,R)$:
\begin{enumerate}[label=\textup{(\roman*)}]
\item $R=96/25$, $\rho\le1841/1000$ and $|\beta|\le321/500$;
\item $R=121/40$, $\rho\le143/100$ and $|\beta|\le93/200$;
\item $5/2\le R\le111/40$,
      $\rho\le7R/20+1/3$ and $|\beta|\le3R/8-3/5$.
\end{enumerate}
The structural inequalities \eqref{ms:structural-classical} supply the
moment bounds in \textup{(i)} whenever $C_n\ge\CE$.
\end{proposition}
\begin{proof}
\emph{The initial radius.}

Take $R=96/25$, $r=1841/1000$, $b=321/500$. The structural inequalities~\eqref{ms:structural-classical} give
\[
 \rho<
 \min\left\{\frac{0.3031\cdot0.646}{c_0-0.3031},
              \frac{0.3322\cdot0.429}{c_0-0.3322}\right\}
 =\frac{7361}{4000}<r.
\]
The sign-moment inequality in Section~\ref{ms:moment-domain} gives
$|\beta|<b$, using
\[
 (1+b^2)^2-r^2(1-b^2)=\frac{93288849}{50000000000}>0.
\]
At $u=R<6$,
$H_4=6uA-3(u^2-1)<36A$ and $B<H_4/144<A/4$. Exact arithmetic gives
\[
 -P^\sharp_{r,b}(R)=
 \frac{2214258283536204102723}{10250000000000000000000}>\frac15.
\]
All six ray signs are negative, proving support in $(-R,R)$.

\emph{A first contraction.}
Take $R=121/40$, $r=143/100$, $b=93/200$, which will be useful for the first contraction of the left probability
panels.
Then
\[
 -P_{r,b}(R)=\frac{2533499510174559}{214958080000000000}>\frac1{100}.
\]
All six ray signs are negative.

\emph{The continuous family.} Assume
\[
 \frac52\le R\le\frac{111}{40},\qquad
 \rho\le r_R=\frac7{20}R+\frac13,\qquad
 |\beta|\le b_R=\frac38R-\frac35.
\]
Let $p(R)=-P_{r_R,b_R}(R)$ and $w=R-5/2$. The coefficients of
$p(R)=\sum_{j=0}^5a_jw^j$ are
\[
\begin{array}{c|r}
j&a_j\\ \hline
0&123579925563/1343488000000\\
1&-467155860219/1679360000000\\
2&705959271/13434880000\\
3&8703567/167936000\\
4&-7499223/839680000\\
5&-531441/2099200000
\end{array}
\]
With $h=11/40$, their signs imply throughout $0\le w\le h$
\[
 p'(R)\le a_1+2a_2h+3a_3h^2
       =-\frac{797749885413}{3358720000000}<0.
\]
Consequently
\[
 p(R)\ge p(111/40)
       =\frac{4403736280322409}{214958080000000000}>\frac1{50}.
\]
For rays, freeze each pair $(r_R,b_R)$. The common larger caps
$r_*=131/100$, $b_*=441/1000$ have at $u=5/2$ the six signs
\[
 -\frac{16991}{700000},\quad-\frac{3801}{500},\quad-\frac{1107}{100},
 \quad-\frac{11615533}{7000000},\quad-\frac{706633}{200000},
 \quad-\frac{2057}{1000}.
\]
They majorize the corresponding signs and derivative bounds for every fixed
pair. Combining the negative boundary estimate with these ray bounds proves
\textup{(iii)}. The moment caps are fixed while moving along a ray.
\end{proof}

\begin{theorem}[Elementary support replacement]\label{ms:04:replacement}
Let $n\ge42$ and let $(P,H)$ be an attained maximizing pair with
$C_n\ge\CE$ and $C_{n-1}\le C_n$. Once the corresponding moment bounds
have been proved for $P$, Proposition~\ref{ms:04:three-bounds} gives all
$65$ support contractions in the proposed local exclusion: one initial
radius $96/25$, one contraction to $121/40$, and $63$ applications of the
affine family with $5/2\le R\le111/40$.
\end{theorem}
The structural inequalities~\eqref{ms:structural-classical} give the
first moment bounds. The second application covers the moment maxima
$71458741/50000000$ and $46459359/100000000$. The other $63$ parameter
triples satisfy the affine bounds of Proposition~\ref{ms:04:three-bounds};
they are listed in the accompanying data. Each contraction uses moment
bounds obtained before that contraction is applied.
